\documentclass{article}
\usepackage[utf8]{inputenc}
\usepackage{amsmath}
\usepackage{geometry}
\geometry{margin=2cm}
\begin{document}

\section*{Questão 114 \, --- \, ENEM 2024 \, --- \, Ciências da Natureza}

\subsection*{Tema: Equilíbrio e Centro de Massa em Rodas}

A questão solicita determinar o ponto correto para fixar uma peça de chumbo na roda para corrigir o deslocamento do centro de massa (CM).

\subsection*{Termos-chave}
\begin{description}
  \item[Centro de Massa:] Ponto onde se pode considerar concentrada a massa de um objeto para análise de movimento e equilíbrio.
  \item[Equilíbrio:] Estado em que forças ou pesos aplicados causam estabilidade, sem oscilações.
  \item[Diametralmente Oposto:] Posição exatamente ao lado oposto em relação a um ponto central.
\end{description}

\subsection*{Análise da Questão}
É pedido identificar o ponto na roda onde deve ser fixada a peça de chumbo para compensar o deslocamento do centro de massa, diminuindo as vibrações e equilibrando a roda.

\textbf{Estratégia para resolução:}

Deve-se observar o deslocamento do centro de massa da roda em relação ao seu centro geométrico. O peso adicional deve ser colocado no ponto diametralmente oposto ao deslocamento, para que o centro de massa se aproxime do centro geométrico, corrigindo o desequilíbrio.

Assim, identificar no esquema apresentado o ponto que fica exatamente em oposição ao CM deslocado e confirmar que o peso deve ser fixado ali.

\subsection*{Revisão de conceitos}
\begin{tabular}{|p{4cm}|p{5cm}|p{5cm}|}
\hline
\textbf{Conceito} & \textbf{Ideia-chave} & \textbf{Aplicação na questão} \\
\hline
Centro de Massa (CM) & Localização concentrada da massa total & Identificar posição do CM deslocado na roda para entender desequilíbrio \\
\hline
Equilíbrio de Roda & Contrapeso ao CM deslocado para minimizar vibrações & Peso deve ser colocado no lado oposto ao CM para restaurar equilíbrio \\
\hline
Relação Diametral & Peso deve estar do lado contrário ao deslocamento & Escolha do ponto 3, oposto ao CM deslocado à direita \\
\hline
\end{tabular}

\subsection*{Teoria}

O centro de massa de um objeto é o ponto onde sua massa pode ser considerada concentrada para efeito de análise do movimento e equilíbrio. Em sistemas rotativos, como uma roda com um pneu, o deslocamento do centro de massa em relação ao centro geométrico provoca vibrações indesejadas durante a rotação.

Para corrigir este desequilíbrio, pesos são adicionados em pontos opostos da roda para deslocar o centro de massa para mais próximo do centro geométrico, buscando estabilidade e menor vibração.

Colocar o peso diametralmente oposto ao deslocamento do centro de massa cria um torque que contrabalança o desequilíbrio, estabilizando a roda durante a rotação.

\subsection*{Resolução}

Analisando o esquema, observa-se que o centro de massa (CM) está deslocado para o lado direito da roda, configurando uma concentração maior de massa e causando instabilidade.

Para equilibrar, deve-se colocar o peso adicional no lado oposto, ou seja, à esquerda da roda.

Dentre os pontos indicados para fixar a peça de chumbo, o ponto 3 está exatamente no lado esquerdo, diametralmente oposto ao deslocamento do CM para a direita.

Portanto, a peça deve ser fixada no ponto 3 para corrigir o posicionamento do centro de massa, reduzindo vibrações e promovendo equilíbrio na rotação.

\subsection*{Pulo do Gato}

\textit{Coloque o peso exatamente no lado oposto do deslocamento do centro de massa para garantir o equilíbrio perfeito.}

\subsection*{Análise das alternativas}
\begin{itemize}
  \item[A:] Está do mesmo lado do deslocamento do centro de massa, o que aumentaria o desequilíbrio.
  \item[B:] Embora esteja no lado esquerdo, não é o ponto diametralmente oposto ao centro de massa deslocado, reduzindo a eficiência da correção.
  \item[C:] Ponto diametralmente oposto ao centro de massa deslocado, ideal para fixar peso e corrigir o desequilíbrio.
  \item[D:] Não está alinhado diametralmente ao deslocamento do centro de massa, portanto não corrige adequadamente.
  \item[E:] Fica do mesmo lado do deslocamento do centro de massa, piorando o desequilíbrio.
\end{itemize}

\subsection*{Código da questão}
\texttt{CN\_BALANCEAMENTO\_001}


\end{document}